\documentclass[11pt]{article}
\usepackage[utf8]{inputenc}
\usepackage[T1]{fontenc}
\usepackage{lmodern}
\usepackage[margin=1in]{geometry}
\usepackage{amsmath,amssymb}
\usepackage{textcomp}
\usepackage[dvipsnames]{xcolor}
\usepackage{enumitem}
\usepackage{parskip}
\usepackage{titlesec}
\usepackage[colorlinks=true,linkcolor=blue,citecolor=blue,urlcolor=blue]{hyperref}
\usepackage[round,authoryear]{natbib}

\titleformat{\section}{\large\bfseries}{\thesection}{1em}{}
\titleformat{\subsection}{\normalsize\bfseries}{\thesubsection}{1em}{}
\titleformat{\subsubsection}{\normalsize\bfseries\itshape}{}{0em}{}

\title{\textbf{Enforceable, Not Certifiable:\\ Locating the Boundary Between Protocol-Layer and Model-Layer Negotiation Defense}}
\author{}
\date{}

\newcommand{\ph}{\textit{[placeholder --- fill in after experiments are run]}}

% Camera-ready line-breaking pass (stage_review.md Phase G [minor]): absorb the three
% cosmetic Overfull \hbox warnings without altering any content.
\sloppy

\begin{document}
\maketitle
\vspace{-2em}

\begin{quote}
\small\textit{Assembled from the six confirmed stage outputs of an \texttt{auto-research-discussion} run (seed topic B5.3), post research-taste audit (\textcolor{ForestGreen}{\textbf{PASS}}). Status: \textbf{specification / paper draft --- no experiments have been run.} All numbers below are pre-registered design targets, not results.}
\end{quote}

\begin{abstract}
A protocol can make an action unprofitable only if that action lives inside the space the protocol can observe and mechanically enforce. Negotiation manipulation between autonomous LLM agents does not obviously respect that boundary: some categories --- reneging on a commitment, asserting a checkable fact --- sit inside a space a well-built protocol can see and enforce against, while others --- a misrepresented private valuation, an emotionally coercive appeal --- are content that any well-formed, schema-valid message can still carry. Whether protocol-layer mechanism design actually eliminates negotiation manipulation, or only the fraction of it that happens to be enforceable, has not been measured category by category: classical incentive-compatibility guarantees --- delivered by construction for rational bidders over a formal action space --- offer no evidence either way for an open-ended natural-language negotiator, and the first frameworks that do vary structured-protocol constraints as treatments for LLM negotiators score game-theoretic rationality and allocative efficiency, not per-category manipulation success against hidden ground truth under an incentivized attacker. We instantiate five manipulation categories as bilateral bargaining scenarios with hidden, machine-checkable ground truth, cross them with a four-rung protocol ladder implemented as deterministic environment wrappers, score every trial against precomputed truthful-play surplus rather than a message's structural form, isolate the protocol's effect from attacker competence with a dual prompted/scripted attacker design, validate the scoring rule against human judgment before trusting any verdict, and test whether an activation-probe monitor recovers ground on categories no protocol reaches. We report which manipulation categories are driven to near-zero success by protocol design alone, which persist under every tested protocol, and what that split implies for whether model-layer negotiation defense has an irreplaceable role or folds into mechanism-design engineering.
\end{abstract}

\section{Introduction}

As autonomous LLM agents begin negotiating real transactions on behalf of their principals --- clearing a price, splitting a resource, settling contract terms, often without a human reading every message first --- the default engineering answer to keeping these negotiations honest is a protocol-layer one: constrain what an agent is allowed to say, not what model is saying it. Structured offer schemas, escrow-backed commitments, and machine-checked factual claims all promise the same thing classical mechanism design has long promised for rational bidders: that a well-designed action space can make deception unprofitable regardless of who plays it.

Yet whether that promise actually holds for an LLM negotiator has never been measured category by category, and there is a specific reason to doubt it holds uniformly. A protocol can only make an action unprofitable if that action lives inside the space it can observe and mechanically enforce. Negotiation manipulation does not respect that boundary: reneging on a commitment and asserting a checkable external fact are actions a well-built protocol can see and enforce against, but a misrepresented private valuation or an emotionally coercive appeal is content that any well-formed, schema-valid message can still carry. Whether protocol design eliminates negotiation manipulation, then, is not one question but as many questions as there are manipulation categories --- and answering it category by category is exactly what has not been done.

The nearest existing work approaches this from either side of the boundary without crossing it. Auction-design networks such as AMenuNet deliver dominant-strategy incentive compatibility by construction, but strictly over a formal bid action space in which a bidder's entire strategic footprint is a valuation report inside a fixed schema --- the guarantee has nothing to say about a negotiator who can also persuade, fabricate context, or renege through a natural-language channel the mechanism never observes. On the model side, work on strategic dishonesty shows LLMs can produce deceptive output that fools every tested output-level judge while remaining detectable only from internal activations --- direct evidence that content a protocol cannot see can defeat a check built on top of it, though this line of work never tests whether an interface-level restriction could have prevented the deception in the first place. Between these two poles sit evaluation frameworks like SOTOPIA, which scores rich multi-turn negotiation scenarios but has no hidden ground-truth economic optimum to score manipulation success against, and classical protocol-redesign theory, which eliminates deception within formally specified, discrete-deal domains that have no open-ended natural-language claims to manipulate in the first place. A concurrent 2026 wave has since begun crossing the protocol side of this boundary empirically, and we concede its ground plainly: a mechanism-design framework for verifiable agent-to-agent interaction \citep{arXiv260814613} already evaluates LLM negotiators under unstructured dialogue, structured protocols, and structured protocols with runtime verification --- a protocol ladder as an experimental treatment --- and finds that mechanism-level incentive compatibility does not automatically transfer to LLM-agent behavior, which is exactly this paper's motivating premise demonstrated empirically; a supply-chain bargaining study \citep{arXiv260807538} ablates the free-text channel entirely, leaving only structured numeric offers; and TERMS-Bench \citep{arXiv260513909} scores LLM negotiators continuously against environment-held concealed payoffs. But these works measure rationality, efficiency, or surplus under cooperative or profit-seeking prompts --- none instantiates a manipulation-category taxonomy, none fields an incentivized attacker, and none asks which manipulation categories a protocol rung eliminates versus leaves open. None of these lines of work, that is, crosses a manipulation-category axis with a protocol-constraint axis under an incentivized attacker in the same bilateral negotiation setting --- which is precisely the gap this paper is built to close.

Turning this gap into a measurable question surfaces four difficulties of its own. Manipulation is not one thing: five structurally different attack channels --- reported values, asserted facts, affect, coordination, and timing --- must each be scored on its own terms, because a protocol that closes one channel may leave every other one wide open. Even where a protocol constrains message structure, structure and truthfulness are logically independent properties of a message, so a well-formed field can still carry a false claim, and no format-level check can certify content by itself. Any enforcement mechanism strict enough to matter also raises a bootstrapping question: if the checker verifying a commitment or auditing a claim is itself an LLM-mediated process, the check inherits the same manipulability it is meant to rule out. And because a stricter protocol simultaneously changes what the defense can catch and how hard the attack is to execute, a measured drop in manipulation success is genuinely ambiguous between ``the protocol closed the channel'' and ``the attacker failed to adapt to the new format'' unless those two explanations are pulled apart by design.

Our approach turns each of these obstacles into a specific design choice. Every manipulation category gets its own hidden ground truth and its own verdict under a pre-registered, multiplicity-corrected decision rule, so no aggregate number can hide a channel a protocol failed to close. Every trial is scored against that hidden ground truth directly --- a sampled private valuation, a fact table, a commitment ledger --- never against a message's conformance to a schema, so structure and content are never conflated. Commitment enforcement and fact auditing are implemented as deterministic, non-LLM environment code, so the check itself cannot be gamed the way the underlying negotiation can, and the residual claims no such check can reach are named as an explicit boundary rather than assumed away. And every cell is run under both a free-strategy and a capability-floored scripted attacker, so a measured protocol effect can be attributed to the protocol rather than to the attacker's competence --- a design we validate before trusting it, by auditing the scoring rule itself against human judgment, and extend, for whichever categories resist every tested protocol, into a direct test of whether a candidate model-layer defense recovers the ground protocols cannot reach.

We evaluate this design across bargaining scenarios drawn from real negotiation corpora, comparing manipulation success under an unstructured baseline against three increasingly strict protocol rungs, alongside a scripted non-adaptive attacker control and an automated-judge scoring baseline that isolates what an objective, ground-truth-anchored scorer adds.

In summary, we contribute:
\begin{itemize}[leftmargin=1.4em]
\item An identification and operationalization of the observability boundary underlying protocol-layer negotiation defense: a mechanism can only make an action unprofitable inside the space it can see and enforce, and negotiation manipulation splits across that boundary in a way that has not been measured category by category.
\item An objective, ground-truth-anchored manipulation scorer --- validated against human judgment before being trusted --- that scores hidden private state directly rather than a message's structural conformance, so a protocol cannot appear to ``fix'' manipulation merely by changing how false content is packaged.
\item A design that isolates a protocol's causal effect on manipulation success from attacker competence via a dual prompted/scripted attacker construction, and a direct test of whether an activation-probe monitor recovers ground on exactly the categories that resist every tested protocol, without collapsing legitimate agreement.
\item A pre-registered design for producing an empirical, per-category account of which manipulation types protocol design eliminates and which persist under every tested protocol --- and, once run, for adjudicating what that split implies for whether model-layer negotiation defense has irreplaceable scope or is an engineering problem protocol design alone will eventually close.
\end{itemize}

\section{Problem Definition}

\noindent\textbf{Formal setup.} Let $\mathcal{K} = \{a, b, c, d, e\}$ index five manipulation categories --- misrepresented private valuation ($a$), fabricated external fact or urgency ($b$), emotional or rapport-based coercion ($c$), false commitment or reneging ($d$), and collusive signaling ($e$) --- and let $\mathcal{P} = \{P0, P1, P2, P3\}$ index a protocol-constraint ladder of increasing strictness: $P0$ unstructured free text, $P1$ schema-validated structured offers, $P2 = P1$ plus escrow-backed commitment enforcement, $P3 = P2$ plus fact-table-audited claims. A trial $t$ in cell $(k, p)$ yields a transcript $\tau$, an attacker surplus $u_A$, a defender surplus $u_D$, and hidden scenario parameters $\theta$ --- sampled private valuations, a fact table of externally verifiable claims, and a commitment ledger --- none of which is visible to the opposing agent during play. A truthful-play surplus distribution $S^{\mathrm{tru}}_k$ is precomputed from scripted, non-manipulative agents. A manipulation-success indicator $\mathrm{MSR}_{k,p}$ is scored against $\theta$ directly, never against $\tau$'s conformance to a protocol schema. The object of study is, for every $k \in \mathcal{K}$: does $\mathrm{MSR}_{k,p}$ fall to a near-zero, statistically confirmed level under some tested $p$ (protocol-fixable), does it remain above a nontrivial level even at the strictest tested $p$ (protocol-invariant), or is the evidence ambiguous between the two (undecided) --- and for every protocol-invariant $k$, does a candidate model-layer defense reduce $\mathrm{MSR}_{k,P3}$ without reducing the rate at which agents reach any deal at all.

\noindent\textbf{Why this is hard.} Fixing the category set $\mathcal{K}$ does not fix a single verdict rule: a category-level claim is only meaningful if every category carries its own hidden ground truth and its own statistically corrected decision boundary, because a protocol that eliminates one category's manipulation may leave every other category's success rate untouched --- collapsing the five categories into one aggregate rate would hide exactly this heterogeneity. Even fixing $\mathcal{P}$ does not guarantee that a stricter $p$ constrains $\theta$-verified truthfulness rather than merely $\tau$'s surface form: structural conformance to a schema and the truthfulness of the content inside it are logically independent properties of a message, so $\mathrm{MSR}_{k,p}$ must be computed against $\theta$ directly or a purely cosmetic change in packaging could masquerade as a real reduction in manipulation. Enforcement at $P2$ and $P3$ raises a further difficulty: a check is only as trustworthy as whatever verifies it, so if the checker were itself an LLM-mediated process rather than deterministic environment code, ``enforced'' and ``manipulable'' would collapse into the same category for exactly the claims the protocol is meant to certify. And because moving from $p$ to a stricter $p'$ changes both what the protocol can catch and how difficult the attack is to execute in the new format, an observed change in $\mathrm{MSR}_{k,p}$ is confounded between a genuine protocol effect and an attacker-competence effect unless attack execution is measured separately from attack success under both a free-strategy and a capability-floored attacker.

\section{Method Overview}

At a conceptual level, this is a single measurement pipeline with six cooperating parts, run across a factorial grid of five manipulation categories and four protocol rungs. A category-indexed scenario taxonomy supplies every cell's hidden ground truth and the pre-registered decision rule that turns raw trial outcomes into a per-category verdict. A ground-truth content scorer turns each trial's hidden parameters --- never its message form --- into a manipulation-success reading. Deterministic environment-side enforcement implements exactly the escrow and fact-audit checks the protocol ladder promises, without delegating that check back to an LLM. A dual-attacker confound control runs every cell under both a free-strategy and a capability-floored attacker so a measured protocol effect can be attributed to the protocol rather than to attacker competence. A scorer-validity audit checks the ground-truth scorer itself against human judgment before any verdict built on it is trusted. And a joint defense-recovery test takes whichever categories resist every tested protocol and asks whether a candidate model-layer instrument closes that residual gap without suppressing legitimate agreement.

\section{Methodology}

\subsection{Category-Indexed Scenario Taxonomy}
\textbf{Challenge.} Manipulation is not monolithic --- five structurally different attack channels must each be evaluated on its own terms, and a decision rule spanning many simultaneous per-category comparisons risks manufacturing false verdicts from ordinary noise if left uncorrected.

\textbf{Insight.} A category-level verdict rule with built-in multiplicity control is what prevents a single aggregate success number from masking a real invariant category behind fixed ones, or an uncorrected battery of tests from inventing an invariant verdict that isn't really there.

\textbf{Method.} Each of the five manipulation categories is instantiated as a parameterized scenario family shipping a machine-readable ground-truth file: sampled private valuations, a fact table of externally verifiable claims with known truth values, a commitment ledger, and a truthful-play surplus distribution precomputed from scripted, non-manipulative agents. A category is declared protocol-fixable if its manipulation-success rate falls to at most 5\% under some tested rung and is significantly below the unstructured baseline (two-proportion test, Holm--Bonferroni corrected across the full category-by-rung contrast family); protocol-invariant if its success rate remains at least 10\% under the strictest tested rung with the 95\% confidence interval excluding 5\%; and undecided if neither condition is met.

\subsection{Ground-Truth Content Scorer}
\textbf{Challenge.} A protocol constraint on message structure says nothing, by itself, about whether the content inside a well-formed field is true --- a false valuation typed into a schema-valid field is still false.

\textbf{Insight.} Because structural conformance and content truthfulness are logically independent properties of a message, the only way to avoid conflating them is to score success against the hidden parameters a message's content should have matched, never against the message's form.

\textbf{Method.} Manipulation success in every trial is scored directly against that trial's hidden scenario parameters --- deviation from the sampled private valuation, a fact-table-false assertion, a ledger-recorded commitment later broken --- combined with a surplus-quantile condition (a concluded deal, defender surplus below the truthful-play distribution's 25th percentile, and attacker surplus above its median). No part of this scoring inspects whether the message conforms to the active protocol's schema.

\subsection{Deterministic Environment-Side Enforcement}
\textbf{Challenge.} A ``verifiable'' commitment or an ``audited'' claim is only as trustworthy as whatever performs the verification --- if the checker is itself an LLM-mediated process, it inherits the same manipulability as the negotiation it is meant to constrain.

\textbf{Insight.} Grounding enforcement in deterministic, non-LLM code sidesteps this recursive-trust problem for exactly the claims such code can reach, while making the claims it cannot reach --- intentions, subjective valuations --- an explicit, reported residue rather than a silently assumed absence of risk.

\textbf{Method.} Escrow-backed commitment enforcement auto-executes accepted offers from environment-held balances, making reneging mechanically impossible except through a priced decommitment penalty; fact auditing checks any typed factual assertion against the scenario's fact table and strips unverifiable claims before delivery. Both checks are implemented as deterministic environment code, never as a call to a language model, and both apply only to the machine-checkable claim subset the protocol ladder was designed to reach.

\subsection{Dual-Attacker Confound Control}
\textbf{Challenge.} A drop in manipulation success under a stricter protocol is ambiguous between ``the protocol closed the channel'' and ``the attacker failed to execute the attack in the new format,'' and the two explanations imply opposite conclusions about the protocol's real effect.

\textbf{Insight.} A capability-floored attacker that cannot fail to attempt the attack provides the comparison a free-strategy attacker alone cannot: if a protocol's effect survives when the attacker's competence is held fixed, the effect cannot be attributed to competence.

\textbf{Method.} Every (category, protocol) cell is run under two attacker modes --- a free-strategy prompted attacker and a templated scripted attacker with manipulation moves inserted at fixed turns --- with attack execution (whether the manipulation content actually entered the dialogue) scored separately from attack success. A mixed-effects logistic model regresses success on protocol interacted with category, with random intercepts for scenario and model pair, isolating the protocol's main effect from scenario- and pair-level variance.

\subsection{Scorer Validity Audit}
\textbf{Challenge.} Every downstream category verdict is only as trustworthy as the rule-based success criterion it is computed from, and open-ended negotiation transcripts carry no objective label to check that criterion against.

\textbf{Insight.} Treating the scorer's validity as a checked, falsifiable precondition --- rather than an assumption every other claim silently rests on --- is what keeps a systematically wrong scorer from being discovered only after a full-scale conclusion has already been drawn from it.

\textbf{Method.} Before any scale-up verdict is trusted, the rule-based scorer's success labels are audited against human judgment on a held-out sample of 200 transcripts rated by two annotators, computing Cohen's kappa. The scale-up matrix proceeds only if kappa reaches at least 0.7; disagreement cases are examined against each scenario's valuation-asymmetry level rather than discarded.

\subsection{Joint Defense-Recovery Test}
\textbf{Challenge.} Showing that a manipulation category resists every tested protocol establishes only a negative result --- that no protocol on the tested ladder reaches it --- and says nothing about whether any actual defense can.

\textbf{Insight.} Requiring a candidate defense to reduce manipulation success and preserve the agreement rate jointly is what rules out the trivial ``defense'' of simply making the negotiator refuse to trade, mirroring the same discipline used elsewhere in this design to rule out degenerate signals.

\textbf{Method.} On every category classified protocol-invariant, two model-layer interventions are run at the strictest tested protocol: a defended prompt carrying explicit manipulation-awareness and counterpart-claim skepticism instructions, and an activation-probe monitor trained on contrastive honesty data, flagging trials whose attacker activations indicate deceptive intent. An intervention counts as recovering ground only if it reduces manipulation success by at least 0.10 absolute (Holm-corrected across invariant categories and interventions) while keeping the agreement rate within 0.10 of the undefended baseline.

\section{Experiments}

\subsection{Experimental Setup}

\textbf{Datasets.} CraigslistBargain \citep{CraigslistBargain_He2018} --- real single-issue price-negotiation scenarios; supplies realistic parameters for the misrepresented-valuation and fabricated-fact/urgency categories, and its single-issue structure keeps the pilot's ground-truth computation cheap. Deal-or-No-Deal \citep{DealOrNoDeal_Lewis2017} --- multi-issue division scenarios with private point values; the direct template for hidden-valuation scenarios and the scale-up domain that tests whether category verdicts generalize beyond single-issue price bargaining. CaSiNo \citep{CaSiNo_Chawla2021} --- campsite multi-issue negotiations annotated with real human rapport and persuasion strategies; grounds the emotional-coercion category's scenario templates and the scripted attacker's manipulation moves in observed human persuasion tactics rather than invented ones, which is specifically what the confound-control design needs to make its capability floor credible.

\textbf{Baselines} (sanity-check comparators, not competitors to beat --- this is a measurement study, and no research question here is framed as ``our procedure beats baseline X''). An unstructured ($P0$) protocol arm, identical to every stricter rung except for an absent schema, escrow, or audit, is the necessary control for every rung comparison, and doubles as the check for whether a category has any manipulation threat to fix at all. Scripted-truthful-agent play defines the truthful-play surplus distribution every cell is scored against, not a comparator to beat. An automated-judge scorer modeled on SOTOPIA's \citep{SOTOPIA_mM7VurbA4r} evaluation protocol, applied to the same transcripts, quantifies what an objective, ground-truth-anchored scorer adds beyond an LLM-judge rating. Defended-prompt-only and probe-monitor-only arms, run at the unstructured protocol with no schema, escrow, or audit at all, isolate what model-layer defense contributes on its own.

\textbf{Metrics.} Manipulation-success rate (MSR) and its Wilson 95\% confidence interval, per cell; attack-execution rate (AER), the fraction of trials in which the manipulation content actually entered the dialogue; reneging rate; agreement rate (the fraction of trials reaching any deal, guarding against a protocol ``fixing'' manipulation by destroying trade); detection/refusal rate for the model-layer defense arm; and Cohen's kappa for the scorer-validity audit.

\textbf{Scale, power, and verdict timing.} The pilot's episode count decomposes as: 2 model pairs $\times$ 5 categories $\times$ 4 rungs $\times$ 2 attacker modes $\times$ 30 trials = 2{,}400 manipulation trials, plus approximately 1{,}200 scripted truthful-play episodes (5 categories $\times$ 6 scenario instances $\times$ ${\sim}40$ runs) to estimate each scenario's truthful-play surplus distribution --- approximately 3{,}600 negotiation episodes in total, on the order of $10^{8}$ tokens (a mid-hundreds-of-USD API-cost class at current prices); the scale-up ($\geq$4 model families, 3 domains, 100 trials per cell) is a low-thousands-USD class. Verdict timing is a consequence of the verdict rule itself and is pre-registered accordingly: at the pilot's $n=30$ per cell, a category with observed MSR of 10\% has a Wilson 95\% CI of roughly $[0.035, 0.26]$, which cannot exclude the 5\% invariance bound, and a $1/30$ observed rate similarly cannot confirm the $\leq$5\% fixable bound --- so no per-category verdict is issuable at pilot scale by construction. The pilot is therefore a go/no-go and instrumentation check only (scorer validity via RQ5, refusal rates, wrapper logging); fixable/invariant verdicts are issued exclusively at scale-up. Because even $n=100$ leaves the invariance contrast marginal when true MSR is near 10\% (Wilson CI lower bound ${\approx}0.055$), the strictest-rung cells of every category --- the cells the invariance verdict is computed from --- are run at $n \geq 200$ in the scale-up, and this asymmetric allocation is part of the pre-registration rather than a post hoc top-up.

\textbf{Research questions:}
\begin{enumerate}[leftmargin=1.6em]
\item RQ1 --- Does the per-category verdict rule produce a non-degenerate split between protocol-fixable and protocol-invariant categories?
\item RQ2 --- Do escrow-enforced commitments and content-borne manipulation respond differently to the protocol ladder?
\item RQ3 --- Does the collapse in reneging success reflect mechanical enforcement rather than reduced attacker persuasiveness?
\item RQ4 --- Is a measured protocol effect attributable to the protocol itself rather than to attacker competence?
\item RQ5 --- Does the rule-based success scorer agree with human judgment closely enough to trust?
\item RQ6 --- Does a model-layer defense recover ground on categories that resist every tested protocol?
\end{enumerate}

\subsubsection{RQ1: Does the per-category verdict rule produce a non-degenerate split between protocol-fixable and protocol-invariant categories?}
\textbf{Setup.} Full ablation matrix: pilot stage runs 2 model families $\times$ 5 categories $\times$ 4 protocol rungs $\times$ 6 scenario instances $\times$ 30 trials per cell ($\approx$3,600 negotiation episodes including truthful-play runs; see Scale, power, and verdict timing); scale-up runs $\geq$4 model families (including cross-model pairs) $\times$ 3 domains $\times$ 100 trials per cell, with strictest-rung cells at $n \geq 200$. The per-category verdict rule is applied only to scale-up cells; the pilot is a go/no-go and instrumentation check.

\textbf{Expected outcome.} At least one of the five categories is classified protocol-fixable and at least one is classified protocol-invariant under the pre-registered thresholds --- not every category landing on the same side.

\textbf{Falsification condition.} Refuted if all five categories land on the same side of the decision rule: if every category is protocol-fixable, protocol design fully closes negotiation manipulation on the tested ladder and this direction folds into mechanism-design engineering; if every category is protocol-invariant, the tested ladder is defensively inert regardless of category, which would call the ladder's design into question rather than confirm a category-level boundary.

\textbf{Hypothesis (plain language).} Manipulation is not monolithic with respect to protocol defenses: some categories are structurally closable by mechanism design, and others are not.

\textbf{Metrics.} Per-cell MSR with Wilson 95\% CI; per-category verdict (fixable / invariant / undecided).

\textbf{Result.} \ph

\subsubsection{RQ2: Do escrow-enforced commitments and content-borne manipulation respond differently to the protocol ladder?}
\textbf{Setup.} Within the full matrix, compare the MSR trajectory across protocol rungs specifically for false-commitment/reneging against misrepresented-valuation and emotional-coercion, holding model family and domain fixed.

\textbf{Expected outcome.} Reneging MSR falls to at most 5\% and significantly below the unstructured baseline by the escrow-enforced rung; misrepresented-valuation and emotional-coercion MSR remain at least 10\% at the strictest tested rung with confidence intervals excluding 5\%.

\textbf{Falsification condition.} Refuted if reneging MSR remains at or above 10\% even under escrow enforcement (mechanical enforcement doesn't work as designed), or if misrepresented-valuation or emotional-coercion MSR also falls to 5\% or below under the strictest protocol (content-borne manipulation turns out to be protocol-fixable after all).

\textbf{Hypothesis (plain language).} Escrow enforcement and content truthfulness are logically independent axes: closing one does not close the other.

\textbf{Metrics.} Per-category, per-rung MSR and confidence interval; fixable-versus-invariant category contrast.

\textbf{Result.} \ph

\subsubsection{RQ3: Does the collapse in reneging success reflect mechanical enforcement rather than reduced attacker persuasiveness?}
\textbf{Setup.} For the false-commitment/reneging category specifically, compare the attack-execution rate against the manipulation-success rate across the protocol ladder under the scripted-attacker arm, where attack attempts are guaranteed by template rather than left to the attacker's own strategy.

\textbf{Expected outcome.} Attack-execution rate stays high across the ladder --- the scripted attacker keeps attempting reneging --- even as manipulation success collapses under escrow enforcement, showing the drop is in success, not in attempt.

\textbf{Falsification condition.} Refuted if attack-execution rate itself collapses alongside manipulation success under escrow enforcement, which would mean the ladder is suppressing the attempt rather than mechanically defeating it once attempted --- a materially different mechanism than the one claimed.

\textbf{Hypothesis (plain language).} The protocol ladder mechanically defeats realized reneging attempts; it does not discourage the attempt itself.

\textbf{Metrics.} Attack-execution rate and manipulation-success rate per protocol rung, scripted-attacker arm, false-commitment/reneging category only.

\textbf{Result.} \ph

\subsubsection{RQ4: Is a measured protocol effect attributable to the protocol itself rather than to attacker competence?}
\textbf{Setup.} Compare the sign of each protocol rung's effect on manipulation success between the prompted and scripted attacker arms, conditional on a nonzero attack-execution rate, across all five manipulation categories.

\textbf{Expected outcome.} The protocol effect's sign agrees between the prompted and scripted arms in at least four of the five categories.

\textbf{Falsification condition.} Refuted if the protocol effect's sign disagrees between the two attacker arms in a majority of categories, which would mean apparent protocol effects are largely an artifact of attacker competence rather than a genuine defensive effect of the protocol.

\textbf{Hypothesis (plain language).} Category verdicts reflect what the protocol actually defends against, not how well the attacker adapted to a new message format.

\textbf{Metrics.} Per-category sign concordance between attacker arms; the protocol-by-category interaction term from the mixed-effects model.

\textbf{Result.} \ph

\subsubsection{RQ5: Does the rule-based success scorer agree with human judgment closely enough to trust?}
\textbf{Setup.} Two annotators independently label a held-out sample of 200 transcripts for whether manipulation occurred; compare against the rule-based scorer's labels via Cohen's kappa, and examine the distribution of disagreement cases across the scenarios' valuation-asymmetry quartiles.

\textbf{Expected outcome.} Cohen's kappa reaches at least 0.7, and disagreement cases concentrate in the highest-valuation-asymmetry quartile, where aggressive-but-honest bargaining most resembles exploitation, rather than being spread uniformly.

\textbf{Falsification condition.} Refuted if kappa falls below 0.7 (the scorer is not yet trustworthy as built, and every downstream verdict must be treated as provisional until it is revised), or if disagreements are spread uniformly across asymmetry levels rather than concentrated (meaning the scorer's error is not even a structured, correctable risk).

\textbf{Hypothesis (plain language).} The rule-based scorer is a valid proxy for human judgment of manipulation, with a bounded and structured failure mode rather than arbitrary noise.

\textbf{Metrics.} Cohen's kappa; asymmetry-quartile-conditioned disagreement rate.

\textbf{Result.} \ph

\subsubsection{RQ6: Does a model-layer defense recover ground on categories that resist every tested protocol?}
\textbf{Setup.} On every category classified protocol-invariant by RQ1, run the defended-prompt and activation-probe-monitor interventions at the strictest tested protocol, and compare manipulation success and agreement rate against the undefended arm at that same protocol.

\textbf{Expected outcome.} The activation-probe monitor reduces manipulation success by at least 0.10 absolute ($p<0.05$) while keeping the agreement rate within 0.10 of the undefended baseline; the defended-prompt intervention produces a smaller and less reliable reduction.

\textbf{Falsification condition.} Refuted if neither intervention achieves a 0.10 absolute reduction on any protocol-invariant category, or if a reduction is achieved only alongside an agreement-rate drop larger than 0.10 --- a ``defense-by-not-trading'' outcome that would not count as genuine recovery.

\textbf{Hypothesis (plain language).} Activation-probe monitoring is the more effective of the two tested model-layer defenses at closing the residual gap protocols cannot reach.

\textbf{Metrics.} Manipulation-success-rate reduction (undefended minus defended, at the strictest protocol) per intervention; agreement-rate difference; probe held-out AUC, reported before the probe is used as a defense signal.

\textbf{Result.} \ph

\vspace{0.5em}
\noindent\textbf{Further evaluation.} With more time, a natural extension is to test partner-choice, many-to-many negotiation markets rather than fixed bilateral pairs --- explicitly out of scope here because it changes the unit of analysis from a single dyad's manipulation success to a population-level equilibrium, and would need its own identifiability design rather than an extension of the per-category verdict rule above. A second natural extension is testing protocol rungs beyond the tested ladder (e.g., reputation-weighted claim auditing) on any category this study classifies protocol-invariant, since an invariance verdict is only ever a claim about the tested ladder, not about protocol design in general.

\section{Related Work}

The closest prior work to this paper sits on opposite poles of the boundary this paper measures, and it is worth stating the delta plainly before surveying the wider landscape: auction-design networks such as AMenuNet \citep{AMenuNet_cNb5hkTfGC} already deliver incentive compatibility by construction, work on strategic dishonesty \citep{StrategicDishonesty_IbDr8xgUMW} already shows that model-layer probing can catch deception that fools output-level checks, and a concurrent mechanism-design framework for verifiable agent-to-agent negotiation \citep{arXiv260814613} already varies structured-protocol constraints as an experimental treatment for LLM negotiators --- but none of the three crosses a manipulation-category axis with a protocol-constraint axis under an incentivized attacker in the same bilateral negotiation setting, which is precisely the identifiability gap this paper is built to close.

\textbf{Incentive-compatible mechanism design over formal action spaces.} AMenuNet \citep{AMenuNet_cNb5hkTfGC} learns affine-maximizer auction parameters with a permutation-equivariant network so the resulting mechanism is dominant-strategy incentive-compatible and individually rational by construction, in contrast to learned mechanisms whose incentive compatibility is only approximate. A related line of work, UCB-S \citep{UCBS_lsxeNvYqCj}, designs incentive-aware bandit allocation rules to discipline strategic manipulation of a single scalar observable, such as a content provider's click-through rate, using confidence-based screening as a credible threat. Both deliver strong guarantees, but only over a formal, well-defined action space --- a bidder's valuation report, a provider's click-through rate --- that has nothing to say about a negotiator who can also persuade, fabricate context, or renege through a natural-language channel the mechanism never observes.

\textbf{Model-layer deception detection and evaluation.} Work on strategic dishonesty \citep{StrategicDishonesty_IbDr8xgUMW} shows that models can produce fake harmful compliance that fools every tested output-level judge, and that this deception becomes detectable only by monitoring internal activations with linear probes; a companion line of work, Among Us \citep{AmongUs_XP3v1THxsq}, elicits similar instrumental deception in a text-based social-deduction sandbox and tests whether contrastively-trained activation probes generalize to it out of distribution. A separate structural-causal-games framework \citep{HonestyBestPolicy_EmxpDiPgRu} gives a formal, behaviorally grounded definition of deception with graphical detection criteria. Each of these establishes that model-layer instruments can catch deception a protocol-level check would miss, but none tests whether an interface-level restriction on what a model may emit could have prevented the deceptive behavior in the first place, nor whether the causal-games criteria are operational as a scalable scoring instrument for open-ended negotiation transcripts that don't come with a known causal structure. Two more recent studies come closer to this paper's own defense-arm ablation without crossing a protocol-constraint axis: one tests whether in-context prompting and guardrail models reduce dark-pattern manipulation of web agents and finds both only partially effective, with manipulation rising alongside model capability \citep{arXiv251222894}; another tests whether a model-layer prompt scaffold and per-token-KL distillation can keep an agent loyal to its principal against an adversarial counterparty, finding both interventions move along but never cross a leak/over-refusal trade-off frontier \citep{arXiv260630383}. Neither varies a protocol/mechanism-design axis alongside the defense arm, which is exactly the joint test this paper runs. A further caution comes from a concurrent negative-results study of pre-action activation-probe monitoring, which finds a probe that perfectly separates its own construction-time contrast data crosses its decision threshold on essentially none of a set of real pre-action contexts, and that a probe's steering effect is not concept-specific once compared against unrelated semantic-direction controls \citep{arXiv260630449} --- direct rival evidence that this paper's own activation-probe defense arm must be validated on its own held-out data rather than assumed to transfer from the contrastive dataset it was trained on.

\textbf{Negotiation-evaluation scaffolds and classical protocol theory.} SOTOPIA \citep{SOTOPIA_mM7VurbA4r} situates LLM agents in procedurally generated mixed-motive social scenarios, including negotiation, and scores multi-turn interactions along several behavioral dimensions with an automated-judge framework --- the closest existing evaluation scaffold for goal-driven LLM-vs-LLM negotiation, reused here as a construct-validity baseline rather than a results comparison, since it has no hidden ground-truth economic optimum per scenario to score manipulation success against. Classical work on designing conventions for automated negotiation \citep{RosenscheinZlotkin1994} shows that redesigning the deal-representation protocol itself can eliminate specific deception strategies within formally specified, discrete task-oriented domains --- the direct intellectual ancestor of ``protocol-layer mechanism design fixes manipulation,'' but one that never confronts open-ended natural-language claims, because its domains have none to begin with.

\textbf{Protocol-constraint treatments in LLM negotiation.} A concurrent 2026 line of work has begun varying the protocol layer itself for LLM negotiators, and this paper concedes its ground explicitly. A mechanism-design framework for verifiable multi-agent interaction over A2A/MCP \citep{arXiv260814613} encodes alternating-offers bargaining and VCG-style auctions as constraints over agent-to-agent message schemas, adds a runtime verification-and-repair layer, and evaluates LLM backbones under unstructured dialogue, structured protocols, and structured protocols with verification --- finding that verification reduces outcome variance, that structured protocols raise completion, and, most relevantly here, that mechanism-level incentive compatibility does not automatically transfer to LLM-agent behavior (one model bids its exact valuation in every auction trial; another does so in 3.3\% of trials). That framework is this paper's closest apparatus neighbor: it owns the protocol-ladder-as-treatment design, but it measures deviations from game-theoretic predictions under non-adversarial prompts --- no manipulation-category taxonomy, no incentivized attacker, no hidden-ground-truth manipulation scoring, no per-category fixable/invariant verdict, and no model-layer defense arm. A supply-chain bargaining study \citep{arXiv260807538} similarly ablates the verbal channel entirely --- agents exchange only numeric offers --- as one contrast on surplus division under private information, and an information-regime study of honesty and credulity in simulated bargaining \citep{arXiv260531445} finds off-the-shelf LLMs attempt to lie about private information but exploit asymmetries poorly, with profit-maximizing fine-tuning making agents both stronger and more dishonest --- direct evidence for the misreport threat this paper's category ($a$) instantiates, measured across information regimes rather than protocol constraints. TERMS-Bench \citep{arXiv260513909}, finally, supplies the nearest scoring apparatus: environment-as-verifier diagnostics against concealed latent types and payoffs, under one fixed protocol --- continuous surplus diagnostics without a manipulation treatment or a protocol axis. Relative to all four, this paper's contribution is the cross itself: per-category manipulation-success verdicts across a protocol ladder under an incentivized, competence-controlled attacker, with a defense arm on exactly the categories the ladder cannot close.

\textbf{Persuasion measurement and security attribution methodology.} Work recasting LLM persuasion as Bayesian persuasion \citep{StrategicPersuasion_aTCXvJKnkE} converts debate and persuasion datasets into sender-receiver games with explicit states, priors, and utilities, measuring persuasion as belief-sensitive utility gain rather than rated convincingness --- but within one fixed interaction protocol, never varied as an experimental treatment, so it cannot say which persuasion effects a stricter protocol would remove. Separately, work attributing agent security failures \citep{ThreatSnapshots_kga18ld70t} isolates a single LLM call with its full execution context to attribute failures to the backbone model rather than to surrounding tool code --- the layer-attribution methodology this paper's design borrows in spirit, applied there to single-call prompt-injection attacks rather than to a multi-turn, economically incentivized bargaining interaction where the ``attack'' is distributed across a whole dialogue.

\section{Conclusion}

This work turns an assumption implicit in how agent-to-agent negotiation is currently defended --- that a sufficiently well-designed protocol can make manipulation unprofitable regardless of which model plays a negotiator --- into a measured, per-category claim rather than an article of engineering faith. The central methodological wager is that ``protocol-fixable'' and ``protocol-invariant'' are only trustworthy labels once every category carries its own hidden ground truth, its own multiplicity-corrected decision rule, and its own check that a measured effect reflects the protocol rather than the attacker's competence or the scorer's own blind spots; a design that skipped any one of those checks could report a boundary that isn't really there. Whichever split the ablation matrix produces, the finding is decision-relevant either way: a non-degenerate split would mean model-layer negotiation defense research has an irreplaceable target, and the defense-recovery test would say whether an instrument already exists for it; a degenerate split toward full fixability would mean this problem folds into protocol-standardization engineering after all. Two limitations bound this claim from the outset. First, any protocol-invariant verdict is relative to the tested four-rung ladder and model families --- an untested protocol could close a category labeled invariant here, and the claim should be read as ``invariant under this ladder,'' not ``invariant to any protocol design.'' Second, the residue this study is built to expose --- claims that lie outside any fact table, such as stated intentions or private tastes --- is a structural boundary that no protocol rung on any ladder is likely to close by construction, which is exactly why testing whether a model-layer instrument can reach it, rather than only proving that protocols cannot, is this paper's final and most consequential test.

\vspace{1em}
\hrule
\vspace{1em}

\section*{Pre-handoff checklist}
\begin{itemize}[leftmargin=1.4em]
\item \textbf{Organization/navigation} --- narrative flows problem $\to$ method $\to$ experiments $\to$ related work without forward references to undefined terms.
\item \textbf{Terminology discipline} --- ``manipulation category,'' ``protocol ladder,'' ``manipulation-success rate,'' ``protocol-fixable/protocol-invariant/undecided'' each defined at first use in Problem Definition, used consistently thereafter; no internal stage IDs (\texttt{sp\_*}, \texttt{m\_*}, \texttt{rq\_*}, C1--C5, M1--M6) leaked into prose --- checked directly against the discussion pipeline log.
\item \textbf{Notation with worked example} --- the formal setup section defines $(\mathcal{K}, \mathcal{P}, \theta, \tau, u_A, u_D,$ $S^{\mathrm{tru}}_k, \mathrm{MSR}_{k,p})$ and the Problem Definition/Methodology sections consistently reuse this notation; no separate worked numeric example is included since no numbers exist yet (would require fabricating results).
\item \textbf{RQ framing matches outcomes} --- all 6 RQs are graded/quantitative (rates, confidence intervals, concordance proportions), each answered in its own subsection with an explicit falsifiable expected outcome and a distinct falsification condition, not a binary yes/no.
\item \textbf{No typesetting artifacts} --- only explicit placeholder markers in Results, per the drafting convention this is a pre-registered spec, not a report.
\item \textbf{Artifact plan} --- the scenario-generator/protocol-wrapper harness and the self-generated transcript corpus are the natural release artifacts; not yet built, so not claimed as available.
\end{itemize}

\section*{Hand-off}
This is a draft, not a finished paper --- experiments have not run. Natural next steps: prose polish and venue-specific LaTeX formatting once a target venue is chosen, or revising an earlier stage if piloting surfaces a problem (e.g., if the scorer-validity kappa in RQ5 is much lower than expected).

\bibliographystyle{plainnat}
\bibliography{refs}

\end{document}
